\clearpage
\section*{四、实验心得与思考题}

\subsection*{1. 实验心得}

通过本次实验，我对并联电容提高功率因数的过程有了更直观的认识。
\begin{enumerate}
\item 随着电容增大，功率因数先升高后降低，总电流则在最佳档位附近达到最小值，说明电容补偿并非越大越好。
\item 同时，同一频率下的有功功率变化较小，让我理解了补偿主要是减小无功电流和线路损耗，而不是降低灯管本身消耗的功率。
\item 比较三种频率的数据后，我发现频率越高，最佳档位对应的电容越小，因此选择补偿电容时也要考虑电源频率。
\item 最后，强电实验必须重视操作顺序，接线和更换电容前先断电，确认电容安全放电后再操作。

\end{enumerate}
\subsection*{2. 思考题}

\subsubsection*{1. 提高功率因数的意义何在？为什么并联电容能提高功率因数？电容是否越大越好？}

\begin{enumerate}
\item 在电压和有功功率一定时，提高功率因数可减小线路电流，降低导线的 $I^2R$ 损耗和电压降，提高供电设备的利用率。
\item 感性负载电流滞后于电压，电容电流超前于电压。两者的无功分量相互抵消，使总电流与电压的相位差减小。
\item 电容不是越大越好。超过适当补偿值后，电路出现过补偿，总电流增大，功率因数反而下降。

\end{enumerate}
\subsubsection*{2. 随电容增加，总电流为什么先减小后增大？}

在正弦交流近似下，将负载电流分为有功分量 $I_P$ 和感性无功分量 $I_Q$，则并联电容后的总电流为

\[
I=\sqrt{I_P^2+(I_Q-2\pi fCU)^2}.
\]

\begin{enumerate}
\item 欠补偿时，增加电容使剩余无功电流减小，总电流随之下降。
\item 接近完全补偿时，两种无功电流基本抵消，总电流达到最小值。
\item 继续增加电容，容性无功电流占主导，总电流重新增大。本次50 Hz和55 Hz的测量清楚地体现了这一趋势。

\end{enumerate}
\subsubsection*{3. 为什么采用并联电容，而不采用串联电容？}

\begin{enumerate}
\item 并联电容补偿无功电流，在电源电压不变时，基本不影响原日光灯支路的工作状态。
\item 串联电容会改变支路阻抗和电压分配，影响灯管工作；接近谐振时还可能产生过电流或过电压。因此，本实验采用并联补偿。
\end{enumerate}
